% F17：本书原创矢量示意；依据所在节的定义、证明或明确模型。
\begin{figure}[htbp]
\centering
\begin{tikzpicture}[x=1cm,y=1cm,>=stealth,every node/.style={font=\small},line width=.65pt]
\draw[->,gray] (-0.12,0)--(5.5,0) node[right] {\(x\)};\draw[->,gray] (0,-0.12)--(0,2.8);
\draw[->,gray] (6.88,0)--(12.5,0) node[right] {\(x\)};\draw[->,gray] (7,-0.12)--(7,2.8);
\draw[AnalysisBlue,thick] plot coordinates {(0.1000,1.5966) (0.1274,1.6160) (0.1547,1.6323) (0.1821,1.6458) (0.2095,1.6566) (0.2369,1.6648) (0.2642,1.6705) (0.2916,1.6741) (0.3190,1.6757) (0.3464,1.6756) (0.3737,1.6740) (0.4011,1.6712) (0.4285,1.6674) (0.4559,1.6630) (0.4832,1.6582) (0.5106,1.6532) (0.5380,1.6483) (0.5654,1.6438) (0.5927,1.6399) (0.6201,1.6367) (0.6475,1.6344) (0.6749,1.6333) (0.7022,1.6334) (0.7296,1.6348) (0.7570,1.6376) (0.7844,1.6418) (0.8117,1.6475) (0.8391,1.6546) (0.8665,1.6630) (0.8939,1.6727) (0.9212,1.6835) (0.9486,1.6953) (0.9760,1.7080) (1.0034,1.7212) (1.0307,1.7349) (1.0581,1.7487) (1.0855,1.7624) (1.1128,1.7757) (1.1402,1.7883) (1.1676,1.8000) (1.1950,1.8105) (1.2223,1.8195) (1.2497,1.8266) (1.2771,1.8317) (1.3045,1.8344) (1.3318,1.8346) (1.3592,1.8319) (1.3866,1.8263) (1.4140,1.8175) (1.4413,1.8054) (1.4687,1.7899) (1.4961,1.7708) (1.5235,1.7483) (1.5508,1.7221) (1.5782,1.6925) (1.6056,1.6594) (1.6330,1.6229) (1.6603,1.5832) (1.6877,1.5405) (1.7151,1.4949) (1.7425,1.4467) (1.7698,1.3962) (1.7972,1.3437) (1.8246,1.2895) (1.8520,1.2339) (1.8793,1.1773) (1.9067,1.1201) (1.9341,1.0627) (1.9615,1.0055) (1.9888,0.9489) (2.0162,0.8932) (2.0436,0.8389) (2.0709,0.7863) (2.0983,0.7358) (2.1257,0.6879) (2.1531,0.6427) (2.1804,0.6007) (2.2078,0.5620) (2.2352,0.5270) (2.2626,0.4959) (2.2899,0.4688) (2.3173,0.4459) (2.3447,0.4273) (2.3721,0.4132) (2.3994,0.4034) (2.4268,0.3980) (2.4542,0.3970) (2.4816,0.4003) (2.5089,0.4077) (2.5363,0.4191) (2.5637,0.4344) (2.5911,0.4532) (2.6184,0.4754) (2.6458,0.5006) (2.6732,0.5286) (2.7006,0.5591) (2.7279,0.5917) (2.7553,0.6262) (2.7827,0.6620) (2.8101,0.6990) (2.8374,0.7367) (2.8648,0.7749) (2.8922,0.8131) (2.9196,0.8511) (2.9469,0.8886) (2.9743,0.9252) (3.0017,0.9608) (3.0291,0.9950) (3.0564,1.0278) (3.0838,1.0587) (3.1112,1.0879) (3.1385,1.1150) (3.1659,1.1401) (3.1933,1.1631) (3.2207,1.1840) (3.2480,1.2028) (3.2754,1.2195) (3.3028,1.2342) (3.3302,1.2470) (3.3575,1.2581) (3.3849,1.2677) (3.4123,1.2758) (3.4397,1.2828) (3.4670,1.2888) (3.4944,1.2941) (3.5218,1.2989) (3.5492,1.3035) (3.5765,1.3080) (3.6039,1.3129) (3.6313,1.3182) (3.6587,1.3243) (3.6860,1.3313) (3.7134,1.3395) (3.7408,1.3491) (3.7682,1.3601) (3.7955,1.3728) (3.8229,1.3873) (3.8503,1.4037) (3.8777,1.4219) (3.9050,1.4421) (3.9324,1.4641) (3.9598,1.4881) (3.9872,1.5138) (4.0145,1.5412) (4.0419,1.5702) (4.0693,1.6005) (4.0966,1.6319) (4.1240,1.6643) (4.1514,1.6974) (4.1788,1.7309) (4.2061,1.7645) (4.2335,1.7979) (4.2609,1.8307) (4.2883,1.8627) (4.3156,1.8936) (4.3430,1.9229) (4.3704,1.9503) (4.3978,1.9756) (4.4251,1.9983) (4.4525,2.0183) (4.4799,2.0352) (4.5073,2.0488) (4.5346,2.0588) (4.5620,2.0650) (4.5894,2.0673) (4.6168,2.0655) (4.6441,2.0595) (4.6715,2.0492) (4.6989,2.0346) (4.7263,2.0158) (4.7536,1.9927) (4.7810,1.9654) (4.8084,1.9341) (4.8358,1.8990) (4.8631,1.8601) (4.8905,1.8179) (4.9179,1.7724) (4.9453,1.7240) (4.9726,1.6731) (5.0000,1.6199)};
\node[] at (2.5,3.3) {原函数};
\node[] at (9.5,3.3) {单元平均};
\draw[gray!50](7.0,0)--(7.0,2.4);\draw[AnalysisBlue,thick](7.0,1.6363503598293876)--(7.5,1.6363503598293876);\draw[gray!50](7.5,0)--(7.5,2.4);\draw[AnalysisBlue,thick](7.5,1.6550938885601825)--(8.0,1.6550938885601825);\draw[gray!50](8.0,0)--(8.0,2.4);\draw[AnalysisBlue,thick](8.0,1.7979760626180894)--(8.5,1.7979760626180894);\draw[gray!50](8.5,0)--(8.5,2.4);\draw[AnalysisBlue,thick](8.5,1.4027509534013876)--(9.0,1.4027509534013876);\draw[gray!50](9.0,0)--(9.0,2.4);\draw[AnalysisBlue,thick](9.0,0.5639346427607738)--(9.5,0.5639346427607738);\draw[gray!50](9.5,0)--(9.5,2.4);\draw[AnalysisBlue,thick](9.5,0.640714777641407)--(10.0,0.640714777641407);\draw[gray!50](10.0,0)--(10.0,2.4);\draw[AnalysisBlue,thick](10.0,1.1766643798831886)--(10.5,1.1766643798831886);\draw[gray!50](10.5,0)--(10.5,2.4);\draw[AnalysisBlue,thick](10.5,1.3733195895322858)--(11.0,1.3733195895322858);\draw[gray!50](11.0,0)--(11.0,2.4);\draw[AnalysisBlue,thick](11.0,1.8071584232090536)--(11.5,1.8071584232090536);\draw[gray!50](11.5,0)--(11.5,2.4);\draw[AnalysisBlue,thick](11.5,1.9390632896824511)--(12.0,1.9390632896824511);
\end{tikzpicture}
\caption[网格平均与有限维逼近]{网格平均与有限维逼近。有限区间上的网格平均示意。一般 \(L^p\) 论证先控制空间尾部，再以平移估计控制网格平均误差；最后保留的网格数有限。}
\label{b1:fig:visual-f17}
\end{figure}
